\section{Matices Jurídicos Canarios, Salidas Legales y Efecto Escala}
\label{sec:realidad_juridica}

\subsection{Salidas Reales para la Edificabilidad Retenida: TAU, Subsuelo y Bajo Cubierta}
\label{sub:salidas_tau}

Ante los déficits estructurales detectados en los escenarios $S_1$ ($11.400\text{ m}^2\text{c}$) y $S_2$ ($33.000\text{ m}^2\text{c}$) de la Subsección~\ref{sub:escenarios_s1_s4}, resulta imprescindible distinguir las verdaderas vías de aprovechamiento jurídico respecto de los elementos meramente comerciales:

\begin{enumerate}
  \item \textbf{Transferencia de Aprovechamiento Urbanístico (TAU):} Es la única vía ortodoxa de equidistribución para materializar el techo excedente: se edifican en dos plantas los $48.600\text{ m}^2\text{c}$ ($S_1$) o $27.000\text{ m}^2\text{c}$ ($S_2$) que caben físicamente en el sector y se transfiere el aprovechamiento sobrante (mediante coeficientes de ponderación) hacia otro sector receptor previsto por el Plan General con mayor altura permitida.
  \item \textbf{Garajes obligatorios no computables (\emph{sótanos}):} Si las ordenanzas municipales eximen íntegramente del cómputo de edificabilidad las plazas de garaje obligatorias (mínimo legal de $1$ plaza por vivienda según el art.~138 de la Ley 4/2017 \cite{ley4_2017_canarias}), \emph{cualquier} promotor puede construirlas aunque carezca de edificabilidad sobrante. Por ello, los sótanos no computables \textbf{no constituyen una salida matemática} para el techo atrapado, aunque sí mejoran la viabilidad comercial al liberar el $100\%$ de las dos plantas sobre rasante para uso residencial.
  \item \textbf{Supuestos en que el subsuelo o la cubierta SÍ consumen edificabilidad sobrante:}
  \begin{itemize}
    \item \emph{Plazas excedentarias y semisótanos:} Bajo ordenanzas estrictas que solo eximen la dotación mínima obligatoria, una segunda plaza de aparcamiento en unifamiliar, una bodega o un \emph{semisótano} abierto en parcela en pendiente computan parcialmente como techo edificable y absorben parte del déficit.
    \item \emph{Espacios bajo cubierta inclinada:} Un aprovechamiento \emph{bajo cubierta} dentro de la pendiente máxima autorizada ($\sim 30^\circ$--$35^\circ$) y sin superar la altura de cornisa no suele computar como tercera planta de fachada, pero \textbf{sí consume edificabilidad} en toda la zona con altura libre superior a $\sim 1,50$--$1,80\text{ m}$, permitiendo absorber los $0,47$ niveles sobrantes de $S_1$ ($11.400\text{ m}^2\text{c}$) allí donde el planeamiento admita cubiertas inclinadas.
  \end{itemize}
\end{enumerate}

\subsection{La Cesión Lucrativa Municipal del $10\%$--$15\%$ (PMS)}
\label{sub:cesion_municipal}

De acuerdo con el artículo~44.c de la Ley 4/2017 del Suelo de Canarias \cite{ley4_2017_canarias} (cesión del $10\%$, dentro de la horquilla estatal del $5\%$ al $15\%$ establecida por el art.~18.1.b del TRLSRU \cite{trlsru2015}):

\begin{itemize}
  \item \textbf{Objeto de la entrega:} Se ceden al \emph{Patrimonio Municipal de Suelo} (PMS) \textbf{parcelas netas ya urbanizadas, libres de cargas de urbanización}, nunca edificaciones terminadas.
  \item \textbf{Invariancia geométrica del sector en la Tarea 4:} Al quedar integradas las parcelas municipales dentro del $S_{\text{neta}}$ total, el balance global de cesiones EL/D/I y el número mínimo de plantas $N_p$ de la Sección~\ref{sec:modelo_normativo} no se alteran (a diferencia de la \emph{Tarea 3}, que aislaba expresamente el número neto de parcelas atribuibles a los promotores).
  \item \textbf{Asimetría financiera del promotor:} El promotor percibe únicamente el $90\%$ del techo lucrativo ($0,90 \cdot E_{\text{teórica}}$), pero financia el $100\%$ de las obras de urbanización del sector ($S_{\text{bruta}}$).
\end{itemize}

\begin{tcolorbox}[jumpbox]
\textbf{Salto Lateral $\to$ Integración del PMS en la Reserva de VPO:} En sectores mixtos con vivienda protegida, el $10\%$ de cesión municipal computa dentro de la reserva global del $40\%$ de techo protegido, tal como se modela en la Subsección~\ref{sub:desacoplamiento_vpo}.
\end{tcolorbox}

\subsection{Variables Intensivas frente a Escala Extensiva ($S_{\text{bruta}}$)}
\label{sub:escala_extensiva}

Al formularse en términos de cocientes intensivos ($\text{BCR}$ y $\text{Gross FAR}$), las ecuaciones base arrojan idénticos porcentajes para una actuación de $0,5\text{ ha}$ que para un megadesarrollo de $500\text{ ha}$. En la práctica urbanística e inmobiliaria, esa invariancia de escala se rompe por cuatro factores extensivos:
\begin{enumerate}
  \item \textbf{Saltos discretos de parcelación en ámbitos pequeños:} En sectores reducidos ($0,5$--$1,5\text{ ha}$), un pequeño remanente inferior a $S_{\mathrm{min}} = 700\text{ m}^2\text{s}$ provoca pérdidas porcentuales bruscas de unidades edificables.
  \item \textbf{Absorción de mercado y costes financieros:} Saturar un mercado local con cientos de viviendas simultáneas deprime precios y prolonga el periodo de comercialización y financiación.
  \item \textbf{Cargas fijas frente a umbrales normativos:} Las actuaciones pequeñas soportan costes fijos elevados de acometida, mientras que los grandes sectores activan exigencias de \emph{Sistemas Generales} y reservas legales de VPO.
  \item \textbf{Heterogeneidad parcelaria (\emph{Tarea 3}):} La parcela mínima ($700\text{ m}^2\text{s}$) actúa únicamente como cota inferior; el proyecto de reparcelación permite combinar lotes de diferente tamaño y ocupación dentro del mismo sector.
\end{enumerate}
